%% toc: P2 11 11.1,11.2,11.3
\section{Propositional equivalence laws}
\label{sec:p02-w02-p02-007-equivalence-laws}

The normal forms of Section~\ref{sec:prop-syntax-semantics-04} are reached by
rewriting: an implication is eliminated, a negation is pushed inward, a
double negation is cancelled, a conjunction is distributed over a
disjunction.  Each such step replaces one subformula by a logically
equivalent one.  This section states the equivalence laws behind those steps
as separately labelled results, collects them in
Table~\ref{tab:p02-w02-p02-007-equivalence-laws}, and fixes the rule by which
a law is used inside a larger formula.  Every law is proved in the annex from
the recursive truth clauses alone.

\emph{Inputs.} The formula set \(\mathsf{Fm}_{\mathsf{CPL}}\) of
Section~\ref{sec:prop-syntax-semantics-01}, whose primitive constructors are
exactly \(\bottom,\neg,\land,\lor,\to\), with \(\topv\) the abbreviation of
\(\neg\bottom\); the valuations, their extensions \(\widehat v\) and the
relation \(\equiv\) of Section~\ref{sec:prop-syntax-semantics-02}, so that
\(\alpha\equiv\beta\) means that every valuation gives \(\alpha\) and
\(\beta\) the same truth value; and the replacement theorem
\cref{thm:p02-replacement}.
\emph{Conclusions.} The De Morgan, distribution, absorption, double-negation,
material-implication and contraposition laws
(\cref{thm:p02-w02-p02-007-de-morgan,thm:p02-w02-p02-007-distribution,thm:p02-w02-p02-007-absorption,thm:p02-w02-p02-007-double-negation,thm:p02-w02-p02-007-material-implication,thm:p02-w02-p02-007-contraposition});
the lattice, unit, complement and negated-implication laws that complete the
Boolean rewriting repertoire
(\cref{thm:p02-w02-p02-007-lattice,thm:p02-w02-p02-007-bounds,thm:p02-w02-p02-007-negated-implication});
the validity criterion for clauses (\cref{thm:p02-w02-p02-007-tautological-clause});
and an exponential lower bound showing what equivalence-preserving
conversion to conjunctive normal form can cost
(\cref{thm:p02-w02-p02-007-cnf-size}).

\subsection{Law schemas and their instances}
\label{sec:p02-w02-p02-007-schemas}

\begin{hbdefinition}{def:p02-w02-p02-007-law-schema}{Equivalence law schema}
A \emph{law schema} is a pair of expressions \(\sigma,\tau\), written
\(\sigma\equiv\tau\), built by the formation clauses of
\(\mathsf{Fm}_{\mathsf{CPL}}\) from the metavariables \(\varphi,\psi,\chi\),
the constant \(\bottom\) and the abbreviation \(\topv\).  An
\emph{instance} of the schema assigns to each metavariable a formula of
\(\mathsf{Fm}_{\mathsf{CPL}}\), the same formula at every occurrence of
that metavariable in \(\sigma\) and in \(\tau\) together, and expands every
\(\topv\) to \(\neg\bottom\).  The schema is a \emph{law} iff every instance
is a pair of logically equivalent formulas.  Displays omit outer parentheses
and use the precedence order \(\neg\), \(\land\), \(\lor\), \(\to\) of
Section~\ref{sec:prop-syntax-semantics-01}; each display denotes its fully
parenthesized formula.
\index{equivalence law!propositional}
\end{hbdefinition}

Each law below is proved for arbitrary formulas \(\varphi,\psi,\chi\), so
every instance is covered directly.  The same conclusion also follows from the
single instance in which the metavariables are distinct variables, by the
semantic uniform-substitution theorem \cref{thm:p02-uniform-substitution}.

\subsection{The six required law groups}
\label{sec:p02-w02-p02-007-six-laws}

\begin{hblemma}{thm:p02-w02-p02-007-de-morgan}{proved=annex:p02-w02-p02-007-dm}{De Morgan laws}
For all formulas \(\varphi,\psi\in\mathsf{Fm}_{\mathsf{CPL}}\),
\[
\begin{aligned}
 \neg(\varphi\land\psi)&\equiv\neg\varphi\lor\neg\psi,
   &&\text{(De Morgan for }\land\text{)}\\
 \neg(\varphi\lor\psi)&\equiv\neg\varphi\land\neg\psi.
   &&\text{(De Morgan for }\lor\text{)}
\end{aligned}
\]
\index{De Morgan laws!propositional}
\end{hblemma}

\begin{hblemma}{thm:p02-w02-p02-007-distribution}{proved=annex:p02-w02-p02-007-dist}{Distribution laws}
For all formulas \(\varphi,\psi,\chi\in\mathsf{Fm}_{\mathsf{CPL}}\),
\[
\begin{aligned}
 \varphi\land(\psi\lor\chi)&\equiv(\varphi\land\psi)\lor(\varphi\land\chi),
   &&\text{(}\land\text{ over }\lor\text{, left)}\\
 (\psi\lor\chi)\land\varphi&\equiv(\psi\land\varphi)\lor(\chi\land\varphi),
   &&\text{(}\land\text{ over }\lor\text{, right)}\\
 \varphi\lor(\psi\land\chi)&\equiv(\varphi\lor\psi)\land(\varphi\lor\chi),
   &&\text{(}\lor\text{ over }\land\text{, left)}\\
 (\psi\land\chi)\lor\varphi&\equiv(\psi\lor\varphi)\land(\chi\lor\varphi).
   &&\text{(}\lor\text{ over }\land\text{, right)}
\end{aligned}
\]
\index{distribution laws!propositional}
\end{hblemma}

\begin{hblemma}{thm:p02-w02-p02-007-absorption}{proved=annex:p02-w02-p02-007-abs}{Absorption laws}
For all formulas \(\varphi,\psi\in\mathsf{Fm}_{\mathsf{CPL}}\),
\[
 \varphi\land(\varphi\lor\psi)\equiv\varphi,
 \qquad
 \varphi\lor(\varphi\land\psi)\equiv\varphi.
\]
\index{absorption laws!propositional}
\end{hblemma}

\begin{hblemma}{thm:p02-w02-p02-007-double-negation}{proved=annex:p02-w02-p02-007-dn}{Double-negation law}
For every formula \(\varphi\in\mathsf{Fm}_{\mathsf{CPL}}\),
\(\neg\neg\varphi\equiv\varphi\).
\index{double negation!law of}
\end{hblemma}

\begin{hblemma}{thm:p02-w02-p02-007-material-implication}{proved=annex:p02-w02-p02-007-mi}{Material-implication law}
For all formulas \(\varphi,\psi\in\mathsf{Fm}_{\mathsf{CPL}}\),
\(\varphi\to\psi\equiv\neg\varphi\lor\psi\).
\index{material implication!equivalence law}
\end{hblemma}

\begin{hblemma}{thm:p02-w02-p02-007-contraposition}{proved=annex:p02-w02-p02-007-cp}{Contraposition law}
For all formulas \(\varphi,\psi\in\mathsf{Fm}_{\mathsf{CPL}}\),
\(\varphi\to\psi\equiv\neg\psi\to\neg\varphi\).
\index{contraposition!equivalence law}
\end{hblemma}

\subsection{Further Boolean laws}
\label{sec:p02-w02-p02-007-further-laws}

\begin{hblemma}{thm:p02-w02-p02-007-lattice}{proved=annex:p02-w02-p02-007-lat}{Commutativity, associativity and idempotence}
For all formulas \(\varphi,\psi,\chi\in\mathsf{Fm}_{\mathsf{CPL}}\),
\[
\begin{aligned}
 \varphi\land\psi&\equiv\psi\land\varphi, &
 \varphi\lor\psi&\equiv\psi\lor\varphi,\\
 (\varphi\land\psi)\land\chi&\equiv\varphi\land(\psi\land\chi), &
 (\varphi\lor\psi)\lor\chi&\equiv\varphi\lor(\psi\lor\chi),\\
 \varphi\land\varphi&\equiv\varphi, &
 \varphi\lor\varphi&\equiv\varphi.
\end{aligned}
\]
\end{hblemma}

\begin{hblemma}{thm:p02-w02-p02-007-bounds}{proved=annex:p02-w02-p02-007-bd}{Units, annihilators and complements}
For every formula \(\varphi\in\mathsf{Fm}_{\mathsf{CPL}}\),
\[
\begin{aligned}
 \varphi\land\topv&\equiv\varphi, &
 \varphi\lor\bottom&\equiv\varphi, &&\text{(units)}\\
 \varphi\land\bottom&\equiv\bottom, &
 \varphi\lor\topv&\equiv\topv, &&\text{(annihilators)}\\
 \varphi\land\neg\varphi&\equiv\bottom, &
 \varphi\lor\neg\varphi&\equiv\topv, &&\text{(complements)}\\
 \neg\topv&\equiv\bottom. &&&&
\end{aligned}
\]
Here \(\topv\) is \(\neg\bottom\), so the last line is the instance
\(\neg\neg\bottom\equiv\bottom\) of the double-negation law.
\end{hblemma}

\begin{hblemma}{thm:p02-w02-p02-007-negated-implication}{proved=annex:p02-w02-p02-007-ni}{Negated-implication law}
Let \(\varphi\) and \(\psi\) be formulas of \(\mathsf{Fm}_{\mathsf{CPL}}\).
Then
\[
 \neg(\varphi\to\psi)\equiv\varphi\land\neg\psi.
\]
\end{hblemma}

\begingroup\small
\begin{longtable}{@{}p{0.2\textwidth}p{0.5\textwidth}p{0.22\textwidth}@{}}
\caption{Propositional equivalence laws for \(\mathsf{Fm}_{\mathsf{CPL}}\).
Every row holds for all formulas \(\varphi,\psi,\chi\); \(\topv\) abbreviates
\(\neg\bottom\).}
\label{tab:p02-w02-p02-007-equivalence-laws}\\
\toprule Law & Schemas & Result \\
\midrule\endfirsthead
\toprule Law & Schemas & Result \\
\midrule\endhead
De Morgan &
 \(\neg(\varphi\land\psi)\equiv\neg\varphi\lor\neg\psi\);
 \(\neg(\varphi\lor\psi)\equiv\neg\varphi\land\neg\psi\) &
 \cref{thm:p02-w02-p02-007-de-morgan} \\
Distribution &
 \(\varphi\land(\psi\lor\chi)\equiv(\varphi\land\psi)\lor(\varphi\land\chi)\);
 \(\varphi\lor(\psi\land\chi)\equiv(\varphi\lor\psi)\land(\varphi\lor\chi)\);
 and the right-hand forms &
 \cref{thm:p02-w02-p02-007-distribution} \\
Absorption &
 \(\varphi\land(\varphi\lor\psi)\equiv\varphi\);
 \(\varphi\lor(\varphi\land\psi)\equiv\varphi\) &
 \cref{thm:p02-w02-p02-007-absorption} \\
Double negation &
 \(\neg\neg\varphi\equiv\varphi\) &
 \cref{thm:p02-w02-p02-007-double-negation} \\
Material implication &
 \(\varphi\to\psi\equiv\neg\varphi\lor\psi\) &
 \cref{thm:p02-w02-p02-007-material-implication} \\
Contraposition &
 \(\varphi\to\psi\equiv\neg\psi\to\neg\varphi\) &
 \cref{thm:p02-w02-p02-007-contraposition} \\
Commutativity &
 \(\varphi\land\psi\equiv\psi\land\varphi\);
 \(\varphi\lor\psi\equiv\psi\lor\varphi\) &
 \cref{thm:p02-w02-p02-007-lattice} \\
Associativity &
 \((\varphi\land\psi)\land\chi\equiv\varphi\land(\psi\land\chi)\);
 \((\varphi\lor\psi)\lor\chi\equiv\varphi\lor(\psi\lor\chi)\) &
 \cref{thm:p02-w02-p02-007-lattice} \\
Idempotence &
 \(\varphi\land\varphi\equiv\varphi\); \(\varphi\lor\varphi\equiv\varphi\) &
 \cref{thm:p02-w02-p02-007-lattice} \\
Units &
 \(\varphi\land\topv\equiv\varphi\); \(\varphi\lor\bottom\equiv\varphi\) &
 \cref{thm:p02-w02-p02-007-bounds} \\
Annihilators &
 \(\varphi\land\bottom\equiv\bottom\); \(\varphi\lor\topv\equiv\topv\) &
 \cref{thm:p02-w02-p02-007-bounds} \\
Complements &
 \(\varphi\land\neg\varphi\equiv\bottom\);
 \(\varphi\lor\neg\varphi\equiv\topv\); \(\neg\topv\equiv\bottom\) &
 \cref{thm:p02-w02-p02-007-bounds} \\
Negated implication &
 \(\neg(\varphi\to\psi)\equiv\varphi\land\neg\psi\) &
 \cref{thm:p02-w02-p02-007-negated-implication} \\
\bottomrule
\end{longtable}
\endgroup

\subsection{Using a law inside a formula}
\label{sec:p02-w02-p02-007-application}

\begin{hbmethod}{def:p02-w02-p02-007-apply-law}{Rewriting by an equivalence law}
\emph{Input:} a formula \(\varphi\), a law \(\sigma\equiv\tau\) from
Table~\ref{tab:p02-w02-p02-007-equivalence-laws}, and a direction of use.
\begin{enumerate}
  \item Write \(\varphi\) fully parenthesized, so that its parse is the
        unique one of \cref{thm:prop-unique-readability}.
  \item Choose an instance \(\alpha\equiv\beta\) of the law, read in the
        chosen direction: \(\alpha\) is the instance of \(\sigma\) and
        \(\beta\) that of \(\tau\), or the reverse.  Since \(\equiv\) is
        symmetric (\cref{thm:p02-replacement}), both directions are allowed.
  \item Choose a set \(O\) of occurrences of \(\alpha\) as a subformula of
        \(\varphi\), in the sense of \cref{def:p02-selected-replacement}.
  \item Conclude \(\varphi\equiv\varphi[O\leftarrow\beta]\).
\end{enumerate}
\emph{Side conditions.}  (a) The instance must assign one formula to each
metavariable throughout both sides of the law.  (b) Each member of \(O\)
must be a subformula occurrence of the parse in step~1; a string that merely
appears inside the abbreviated display is not enough.  (c) Laws that
mention \(\topv\) apply only to occurrences of \(\neg\bottom\), and a
defined biconditional is expanded to
\((\varphi\to\psi)\land(\psi\to\varphi)\) before any law is applied.
\emph{Correctness.} Step~4 is the replacement clause of
\cref{thm:p02-replacement} applied to the equivalence \(\alpha\equiv\beta\)
given by the cited law.  A chain of such steps yields an equivalence between
its first and last formulas because \(\equiv\) is transitive
(\cref{thm:p02-replacement}).
\index{replacement!by equivalence laws}
\end{hbmethod}

The proof of the negation-normal-form theorem in
Section~\ref{sec:prop-syntax-semantics-04} applies this method: its
elimination of \(\to\) is the material-implication law
(\cref{thm:p02-w02-p02-007-material-implication}), its two inward moves of a
negation are the De Morgan laws (\cref{thm:p02-w02-p02-007-de-morgan}), and
its removal of a double negation is the double-negation law
(\cref{thm:p02-w02-p02-007-double-negation}), each used from left to right at
a selected subformula occurrence.

\begin{hbexample}{def:p02-w02-p02-007-cnf-example}{From an implication to CNF and DNF}
Take \(P=p_0\), \(Q=p_1\), \(R=p_2\) and
\(\varphi=\neg(P\to(Q\land\neg R))\).  Each line below applies one law at
the displayed subformula occurrence, by
\cref{def:p02-w02-p02-007-apply-law}:
\[
\begin{aligned}
 \neg(P\to(Q\land\neg R))
   &\equiv\neg(\neg P\lor(Q\land\neg R))
   &&\text{material implication}\\
   &\equiv\neg\neg P\land\neg(Q\land\neg R)
   &&\text{De Morgan for }\lor\\
   &\equiv P\land\neg(Q\land\neg R)
   &&\text{double negation}\\
   &\equiv P\land(\neg Q\lor\neg\neg R)
   &&\text{De Morgan for }\land\\
   &\equiv P\land(\neg Q\lor R)
   &&\text{double negation.}
\end{aligned}
\]
The last formula is \(p_0\land(\neg p_1\lor p_2)\).  In the coding of
Section~\ref{sec:prop-syntax-semantics-04} the literal \(\neg p_1\) has code
\(3\) and \(p_2\) has code \(4\), so \(\neg p_1\lor p_2\) is exactly the
formula interpretation of the clause \(\{\neg p_1,p_2\}\), and \(p_0\) is
that of \(\{p_0\}\).  Hence the last formula is the conjunctive normal form
with clause list \(\{p_0\},\{\neg p_1,p_2\}\).  One more step, the left
distribution of \(\land\) over \(\lor\), gives
\[
 P\land(\neg Q\lor R)\equiv(P\land\neg Q)\lor(P\land R),
\]
a disjunctive normal form.  As a check, \(\varphi\) is false exactly when
\(P\to(Q\land\neg R)\) is true, that is, when \(P\) is false or both \(Q\)
is true and \(R\) false; the final formula is false in exactly the same
cases.
\end{hbexample}

\begin{hbexample}{def:p02-w02-p02-007-simplification-example}{Contraposition and absorption}
Let \(\varphi=(\neg Q\to\neg P)\land(P\lor(P\land Q))\).  Then
\[
\begin{aligned}
 \varphi
   &\equiv(P\to Q)\land(P\lor(P\land Q))
   &&\text{contraposition, right to left}\\
   &\equiv(P\to Q)\land P
   &&\text{absorption}\\
   &\equiv(\neg P\lor Q)\land P
   &&\text{material implication}\\
   &\equiv P\land(\neg P\lor Q)
   &&\text{commutativity}\\
   &\equiv(P\land\neg P)\lor(P\land Q)
   &&\text{distribution, left}\\
   &\equiv\bottom\lor(P\land Q)
   &&\text{complements}\\
   &\equiv(P\land Q)\lor\bottom
   &&\text{commutativity}\\
   &\equiv P\land Q
   &&\text{units.}
\end{aligned}
\]
The first step uses the instance
\(P\to Q\equiv\neg Q\to\neg P\) of \cref{thm:p02-w02-p02-007-contraposition}
from right to left.
\end{hbexample}

\subsection{Clauses with complementary literals}
\label{sec:p02-w02-p02-007-clauses}

The resolution section (Section~\ref{sec:prop-metatheory-04}) calls a clause
\emph{tautological} when it contains a literal and its complement.  The
complement law \(\varphi\lor\neg\varphi\equiv\topv\) is the reason such a
clause is valid, and the converse also holds.

\begin{hblemma}{thm:p02-w02-p02-007-tautological-clause}{proved=annex:p02-w02-p02-007-taut}{Validity of clauses}
Let \(C\) be a clause and let \(\operatorname{form}_{\mathsf{CPL}}(C)\) be its formula
interpretation from Section~\ref{sec:prop-syntax-semantics-04}.  Then
\(\entails\operatorname{form}_{\mathsf{CPL}}(C)\) iff there is a variable \(p\) with
\(p\in C\) and \(\neg p\in C\).  Consequently, if \(C\) is tautological and
\(S\) is a clause set, a valuation satisfies \(S\) iff it satisfies
\(S\setminus\{C\}\).  The empty clause is not valid.
\end{hblemma}

\subsection{Limitations}
\label{sec:p02-w02-p02-007-limitations}

\begin{hbwarning}{def:p02-w02-p02-007-non-laws}{Look-alike schemas that are not laws}
The converse \(\psi\to\varphi\) and the inverse \(\neg\varphi\to\neg\psi\)
of \(\varphi\to\psi\) are not equivalent to it: with \(\varphi=P\),
\(\psi=Q\) and a valuation making \(P\) false and \(Q\) true,
\(P\to Q\) is true while \(Q\to P\) and \(\neg P\to\neg Q\) are false.
Contraposition reverses \emph{and} negates; doing only one of the two is not
a law.  Likewise \((\varphi\land\psi)\to\chi\) is not equivalent to
\((\varphi\to\chi)\land(\psi\to\chi)\): with \(P\) true, \(Q\) and \(R\)
false, \((P\land Q)\to R\) is true and \((P\to R)\land(Q\to R)\) is false.
Finally, side condition (b) of \cref{def:p02-w02-p02-007-apply-law} is not
a formality: the display \(P\land Q\lor R\) abbreviates
\(((P\land Q)\lor R)\), in which \(Q\lor R\) is not a subformula, so the
distribution law cannot be applied to it as if it read
\(P\land(Q\lor R)\).  The two formulas differ when \(P\) is false and
\(R\) is true.
\end{hbwarning}

Distribution preserves equivalence but not size.  The next theorem shows
that the growth is not an artefact of a careless rewriting order: some short
formulas have no short equivalent conjunctive normal form at all.

\begin{hbtheorem}{thm:p02-w02-p02-007-cnf-size}{proved=annex:p02-w02-p02-007-cnf}{Equivalent conjunctive normal forms can be exponentially long}
Let \(n\geq1\) and let \(P_1,Q_1,\ldots,P_n,Q_n\) be \(2n\) distinct
propositional variables.  Put
\[
 \theta_n\defeq\bigvee_{i=1}^{n}(P_i\land Q_i),
\]
the finite disjunction in the order \(i=1,\ldots,n\); it has \(2n-1\)
binary connective occurrences.  Let \(S_n\) be the set of functions \(s\) on
\(\{1,\ldots,n\}\) with \(s(i)\in\{P_i,Q_i\}\) for every \(i\), and put
\(K_s\defeq\{s(1),\ldots,s(n)\}\).
\begin{enumerate}
  \item The \(2^n\) clauses \(K_s\), \(s\in S_n\), are pairwise distinct, and
        every conjunctive normal form whose clause list enumerates them, each
        once, is logically equivalent to \(\theta_n\).
  \item If a conjunctive normal form is logically equivalent to
        \(\theta_n\), it is not \(\topv\), and its clause list
        \(C_1,\ldots,C_m\) contains at least \(2^n\) distinct clauses; in
        particular \(m\geq2^n\).
\end{enumerate}
\end{hbtheorem}

The definitional conjunctive normal form of
Section~\ref{sec:p02-tseitin-transformation} avoids this growth by adding
fresh variables; it gives up logical equivalence and keeps
equisatisfiability (\cref{thm:p02-tseitin}), with the size bound of
\cref{thm:p02-tseitin-size}.

\subsection{Relationships and sources}
\label{sec:p02-w02-p02-007-relationships}

The laws here concern the semantic relation \(\equiv\).  The Lindenbaum
quotient of the primitive-negation natural deduction satisfies the
corresponding identities between classes of provably equivalent formulas
(\cref{thm:p02-fbc-quotient}); those are statements about derivability,
proved there from the calculus, whereas the lemmas of this section use only
the truth clauses.  The duality theorem (\cref{thm:p02-propositional-duality})
states that dualization preserves logical equivalence; for instances whose
metavariables are negation normal forms, it carries each \(\land\)-law among
distribution, absorption, commutativity, associativity, idempotence, units
and annihilators to its \(\lor\)-partner in
Table~\ref{tab:p02-w02-p02-007-equivalence-laws}, and conversely.  The negation, conjunctive and disjunctive normal forms of
\cref{thm:prop-normal-forms} and Section~\ref{sec:prop-syntax-semantics-04},
and the writing example of Section~\ref{sec:writing-foundations-05}, use the
material-implication, De Morgan and double-negation laws as rewrite steps.

Kleene states the deductive counterparts in his formal propositional
calculus, where ``equivalent'' means that the formula \(A\sim B\), his
abbreviation of \((A\supset B)\mathbin{\&}(B\supset A)\), is provable
\cite[\S 26, p.~113]{Kleene1952}.  In his notation, with \(\&\) for
conjunction and \(\supset\) for implication: the distributive laws are
\({*}35\) and \({*}36\) and the absorption laws are the elimination laws
\({*}39\) and \({*}40\) of Theorem~7
\cite[\S 27, p.~118]{Kleene1952}; the double-negation law is \({*}49\), the
material-implication law is \({*}59\), and the De Morgan laws are
\({*}62\) and \({*}63\), which Kleene attributes to De Morgan in 1847
\cite[\S 27, p.~119]{Kleene1952}.  Kleene gives contraposition as the
deducibility relations \({*}12\) to \({*}15\) of Theorem~5 rather than as an
equivalence \cite[\S 26, p.~113]{Kleene1952}.  He marks \({*}14\), \({*}15\),
\({*}49\), \({*}59\) and \({*}62\) with a degree sign, which signals that the
demonstration he gives is valid for the classical system only
\cite[\S 23, p.~101]{Kleene1952}.
